\hypertarget{diagnostics}{}
\hypertarget{residual-diagnostics-benchmark-sensitivity-and-model-risk}{%
\chapter{Residual diagnostics, benchmark sensitivity and model
risk}\label{residual-diagnostics-benchmark-sensitivity-and-model-risk}}

\hypertarget{diagnostic-tests}{%
\section{Diagnostic tests}\label{diagnostic-tests}}

\begin{longtable}[]{@{}p{0.20\textwidth}p{0.34\textwidth}p{0.36\textwidth}@{}}
\toprule\noalign{}
Test & Null hypothesis & If rejected \\
\midrule\noalign{}
\endhead
\bottomrule\noalign{}
\endlastfoot
Jarque-Bera & Residual skewness and kurtosis match a normal distribution
& Gaussian likelihood is questionable \\
Ljung-Box (lag 12) & No residual autocorrelation through lag 12 & Serial
dependence remains \\
ARCH-LM (lag 12) & No conditional heteroskedasticity & Residual variance
is time-varying \\
\end{longtable}

Each test compares a feature of the residuals to a null and, under that
null, has a $\chi^2$ reference obtained from asymptotically normal
statistics (\S\ref{prelim-tests}).

\textbf{Jarque--Bera (normality).} For residuals $e_1,\dots,e_T$, let $S$ be
the sample skewness and $K$ the sample kurtosis. Under IID normality the
standardised third and fourth moments are asymptotically independent
normals, $\sqrt{T}\,S\xrightarrow{d}N(0,6)$ and
$\sqrt{T}\,(K-3)\xrightarrow{d}N(0,24)$, so their squared standardised
values sum to a chi-square with two degrees of freedom:
\begin{equation}\label{eq:ch13-jb}
JB = \frac{T}{6}\Bigl[S^2 + \tfrac14(K-3)^2\Bigr]
\ \xrightarrow{d}\ \chi^2_2 \quad\text{under an IID normal null}.
\end{equation}

\textbf{Ljung--Box (serial correlation).} Let $\hat\rho_\ell$ be the
residual autocorrelation at lag $\ell$. Under the no-autocorrelation null,
$\sqrt{T}\,\hat\rho_\ell\xrightarrow{d}N(0,1)$ and the $\hat\rho_\ell$ are
asymptotically independent across lags, so a weighted sum of their squares
is chi-square with $m$ degrees of freedom:
\begin{equation}\label{eq:ch13-lb}
Q(m) = T(T+2)\sum_{\ell=1}^{m}\frac{\hat\rho_\ell^{2}}{T-\ell}
\ \xrightarrow{d}\ \chi^2_m .
\end{equation}
The factor $T(T+2)/(T-\ell)$ is the Ljung--Box finite-sample refinement of
the simpler Box--Pierce statistic $T\sum\hat\rho_\ell^2$.

\textbf{ARCH-LM (conditional heteroskedasticity).} This is a
Lagrange-multiplier test built from the auxiliary regression of squared
residuals on their own lags,
\begin{equation}\label{eq:ch13-arch}
\hat e_t^{2} = c + \gamma_1\hat e_{t-1}^{2} + \cdots + \gamma_m\hat e_{t-m}^{2} + \eta_t .
\end{equation}
Under $H_0:\gamma_1=\cdots=\gamma_m=0$ (no ARCH) the auxiliary $R^2$ is
small, and $(T-m)R^2\xrightarrow{d}\chi^2_m$. A large statistic means past
squared residuals predict the current one --- volatility clustering ---
which signals a time-varying conditional variance, not necessarily an
incorrect conditional mean.

\begin{verbatim}
from statsmodels.stats.diagnostic import acorr_ljungbox, het_arch
from statsmodels.stats.stattools import jarque_bera

jb_stat, jb_p, skew, kurt = jarque_bera(resid)
lb_p = acorr_ljungbox(resid, lags=[12], return_df=True)["lb_pvalue"].iloc[0]
arch_stat, arch_p, _, _ = het_arch(resid, nlags=12)
\end{verbatim}

In the paper, normality and no-ARCH are rejected for all 25 portfolios,
while lag-12 independence is rejected for 16. This explains why
IID-normal GRS is secondary, but it does not mean every conclusion is
invalid: HAC and block resampling target specific inferential problems.

\begin{longtable}[]{@{}p{0.29\textwidth}p{0.63\textwidth}@{}}
\toprule\noalign{}
Question & What diagnostics imply \\
\midrule\noalign{}
\endhead
\bottomrule\noalign{}
\endlastfoot
Are OLS coefficients algebraically defined? & Yes, provided X has full
column rank. \\
Are OLS coefficients consistent for the linear projection? &
Potentially, under orthogonality and suitable weak-dependence
conditions; normality is not required. \\
Are conventional IID standard errors valid? & No when heteroskedasticity
or serial dependence is material. \\
Does HAC prove the factor model is correctly specified? & No. It repairs
an asymptotic covariance estimate, not omitted variables, nonlinearities
or unstable betas. \\
Is a Gaussian SFA likelihood trustworthy after these rejections? & Only
as a deliberately narrow diagnostic, reinforced by simulation and
boundary/FDR safeguards. \\
\end{longtable}

\hypertarget{ff3-versus-ff5}{%
\section{FF3 versus FF5}\label{ff3-versus-ff5}}

Alpha is defined only relative to the chosen factors, so enlarging the
benchmark is a genuine source of model risk. FF5 adds profitability (RMW)
and investment (CMA) to FF3. Because each loading is a \emph{partial} slope
(the Frisch--Waugh reading of \S\ref{alpha-sampling}), adding factors
re-partials every coefficient and therefore redefines the intercept: a
portfolio that loads on a priced dimension absent from FF3 shows an FF3
alpha that can shrink or reverse under FF5 (\S\ref{rq-difficulties}). The
two models must be compared on their \emph{common} sample, otherwise a
``model difference'' is confounded with a ``sample difference.''
\begin{itemize}
\tightlist
\item
  Fit both models on the prespecified common sample, July 1963--November
  2025.
\item
  Compare posterior ranks, score correlations, FDR discoveries and the
  joint tests.
\item
  Report the largest absolute rank move and the affected portfolio.
\end{itemize}

\textbf{Paper result.} FF3 and FF5 posterior ranks correlate $0.875$, yet
\texttt{ME5\ BM2} moves from rank $9$ under FF3 to rank $20$ under FF5 --- an
eleven-position change. FDR-supported raw alphas fall from five (FF3) to
three (FF5). The joint HAC/Wald test rejects zero alpha under both models
(FF3 $p=1.6\times10^{-8}$, FF5 $p=6.7\times10^{-7}$), so the \emph{presence}
of mispricing is robust while its \emph{cross-sectional pattern} is not:
benchmark choice is economically meaningful even when the overall rank
correlation looks high.

\hypertarget{window-length-sensitivity}{%
\section{Window-length sensitivity}\label{window-length-sensitivity}}

The rolling window length sets both the sampling noise and the economic
period being averaged, so it is a modelling decision, not a neutral knob.
On aligned dates, ordering agreement between two window lengths is measured
by the Spearman rank correlation and top/bottom-quintile membership
agreement by the Jaccard overlap $J(A,B)=\lvert A\cap B\rvert/\lvert A\cup B\rvert$
(Chapter~10).

\begin{longtable}[]{@{}p{0.22\textwidth}p{0.17\textwidth}p{0.17\textwidth}p{0.17\textwidth}p{0.17\textwidth}@{}}
\toprule\noalign{}
Windows & Rank corr & Score corr & Top Jaccard & Bottom Jaccard \\
\midrule\noalign{}
\endhead
\bottomrule\noalign{}
\endlastfoot
60 vs.\ 120 months & 0.717 & 0.746 & 0.393 & 0.367 \\
60 vs.\ 180 months & 0.546 & 0.580 & 0.325 & 0.294 \\
120 vs.\ 180 months & 0.782 & 0.802 & 0.515 & 0.448 \\
\end{longtable}

Rank correlations are positive but well below one, and top-quintile
membership overlap can fall below $0.40$: two defensible window lengths can
select materially different ``best'' portfolios. This disagreement is
genuine temporal instability, not merely a tuning inconvenience.

\hypertarget{historical-anchor}{%
\section{The historical anchor and the GRS/HAC divergence}\label{historical-anchor}}

A sharp illustration of model risk comes from the prespecified
1963--1991 anchor. There the dependence-robust HAC/Wald test rejects joint
zero alpha ($\mathrm{stat}=56.5$, $p=3.2\times10^{-4}$), while the classical
GRS test does \emph{not} reject at 5\% ($\mathrm{stat}=1.44$, $p=0.082$).
The two disagree precisely because GRS assumes the IID-normal residuals that
the diagnostics of this chapter reject; once those assumptions fail, the
exact-$F$ argument (\S\ref{prelim-tests}) loses its finite-sample guarantee
and can understate the evidence. The divergence is therefore not a
contradiction but a direct demonstration of why the robust test is primary.

\hypertarget{model-risk-summary}{%
\section{Model risk, summarised}\label{model-risk-summary}}

The residual diagnostics, the FF3/FF5 comparison, the window-length
sensitivity and the historical anchor point to one lesson: the linear,
constant-beta, FF3 specification is a useful approximation, not a literal
truth. HAC inference and the block bootstrap make the \emph{inference}
robust to the dependence they model, but they cannot repair omitted priced
factors, nonlinear payoffs, time-varying betas, or an incomplete benchmark.
Every headline claim in the study is therefore stated as
benchmark-relative and paired with the sensitivity analysis that exposes
its model risk rather than concealing it.

